\section{Introduction}

Routine environmental monitoring produces large volumes of physical and chemical
measurements that must be reduced to an interpretable statement about water quality.
Water quality indices serve this purpose by aggregating several parameters into a single
dimensionless score, and the index published by the Canadian Council of Ministers of the
Environment (CCME) is among the most widely adopted~\cite{ccme2001wqi}. Such indices are
attractive targets for supervised learning: they are continuous, widely reported, and
accompanied by exactly the tabular measurement data on which tree ensembles perform
well~\cite{breiman2001randomforests, chen2016xgboost}.

That attraction conceals a methodological trap. A water quality index is not an observed
quantity but a deterministic function of its inputs. When a model is trained on the same
measurements from which the index was computed, it is not learning a relationship between
environmental conditions and water quality; it is approximating a closed-form expression
that is already known. Reported accuracy under this framing is high and the model is
uninformative. In the taxonomy of Kaufman et al.~\cite{kaufman2012leakage} this is target
leakage, and we found it in a published, widely reusable dataset: the distributed CCME
values are reproduced from their eight contributing measurements with zero reconstruction
error.

Our response is to change the question rather than abandon the data. When a monitoring
programme fails to assay a parameter -- because an instrument failed, a sample was lost,
or the assay was never scheduled -- the formula cannot be evaluated at all, yet an index
value is still wanted. Estimating the index from the measurements that \emph{are}
available is a task with genuine content, immediate operational relevance, and no
closed-form solution. We therefore retain the all-measurements setting only as an
emulation benchmark and a check on the pipeline, and define the prediction problem over
subsets of the measurement vector.

Working with a harmonised corpus of 2.8 million monitoring records from England, the USA,
Ireland, China and Canada, we find the distinction is not academic. Withholding a single
measurement moves gradient boosting from $R^2 = 0.9997$ to $0.834$ and multiplies mean
absolute error by a factor of 48. The step between those two numbers is the quantitative
statement of how much apparent accuracy was formula reconstruction rather than inference.

Preparing the data surfaced a second issue of independent interest. The corpus reports no
missing values, but five of its eight measurement columns carry their own exact median as
their most frequent value, at frequencies between 11\% and 32\%. The gaps were filled
before publication, the fill is concentrated almost entirely in one country, and it is
confounded with the target. We treat this as a finding to be measured rather than a
defect to be silently inherited.

This paper contributes the following:
\begin{itemize}
    \item A documented target-leakage pathway in a public water quality dataset, with its
          effect quantified by the gap between full and partial measurement availability.
    \item A partial-measurement evaluation protocol over four regimes, under a
          site-grouped split that admits no monitoring station to more than one partition.
    \item A controlled comparison of Random Forest and XGBoost at every availability
          level, with bootstrap confidence intervals and SHAP attribution.
    \item Detection and quantification of an upstream imputation artefact, together with
          an ablation showing that exposing it as a feature is not justified.
\end{itemize}

The remainder of the paper is organised as follows. Section~\ref{sec:related} reviews
water quality indices, tree ensembles and leakage. Section~\ref{sec:data} describes the
dataset, cleaning and exploratory findings. Section~\ref{sec:method} sets out the
formulation, split and models. Section~\ref{sec:results} reports results,
Section~\ref{sec:limitations} states limitations, and Section~\ref{sec:conclusion}
concludes.
